\documentclass[../../main.tex]{subfiles}

\begin{document}

\section{Tychonoff 定理}\label{section:ch05-tychonoff}

与 Urysohn 引理一样, Tychonoff 定理也是一个``深刻的''定理. 它的证明涉及的思想较新颖, 并非直截了当. 在给出这个定理的证明之前, 我们要详细地讨论在证明中的某些关键性的想法.

第 3 章中, 我们曾证明过两个紧致空间的积 $X\times Y$ 是紧致的. 在那个证明中, 只用紧致性的开覆盖形式就完全够了. 给定 $X\times Y$ 的一个由基元素构成的开覆盖, 先用其中的有限多个元素覆盖每一个薄片 $x\times Y$, 进而构造 $X\times Y$ 的一个有限覆盖.

但对于紧致空间的任意积, 类似于上面的做法已经行不通了, 必须将指标集良序化并且应用超限归纳(参见习题 5). 另外一个办法就是放弃开覆盖方式, 代之以紧致性的闭集形式, 借助 Zorn 引理来处理.

为了弄明白这一思路, 我们先考虑最简单的情形, 即两个紧致空间的积 $X_1\times X_2$ 这种情形. 假定 $\mathcal{A}$ 为 $X_1\times X_2$ 的一个具有有限交性质的闭子集族, 考虑投射 $\pi_1\colon X_1\times X_2\to X_1$. $X_1$ 的子集族
\begin{align*}
\{\pi_1(A)\mid A\in\mathcal{A}\}
\end{align*}
也具有有限交性质, 并且其中元素的闭包 $\overline{\pi_1(A)}$ 的族也具有这个性质. 同时 $X_1$ 的紧致性又保证了所有集合 $\overline{\pi_1(A)}$ 的交非空. 我们从这个交中选择一点 $x_1$. 类似地, 可以从所有集合 $\overline{\pi_2(A)}$ 的交中选择一点 $x_2$. 于是我们自然希望 $x_1\times x_2$ 在 $\bigcap_{A\in\mathcal{A}}A$ 中. 如果果真如此, 定理也就证明了.

但遗憾的是这个设想行不通. 考虑下面的例子: 设 $X_1=X_2=[0,1]$, 集族 $\mathcal{A}$ 是以下所有椭圆形区域族, 其中每一个椭圆形区域的边界都是以 $p=\bigl(\tfrac13,\tfrac13\bigr)$, $q=\bigl(\tfrac12,\tfrac23\bigr)$ 为公共焦点的椭圆, 如\cref{figure:ch05-fig37-1}所示. 当然 $\mathcal{A}$ 具有有限交性质. 现在, 我们在这些集合 $\{\overline{\pi_1(A)}\mid A\in\mathcal{A}\}$ 的交中选取一点 $x_1$, 显然可以选取区间 $\bigl[\tfrac13,\tfrac12\bigr]$ 中的任意一点, 比如选定 $x_1=\tfrac12$. 类似地, 在集合 $\{\overline{\pi_2(A)}\mid A\in\mathcal{A}\}$ 的交中选取一点 $x_2$, 可以选区间 $\bigl[\tfrac13,\tfrac23\bigr]$ 中的任意一点, 比如选定 $x_2=\tfrac12$. 这是一个不幸的选择, 因为点
\begin{align*}
x_1\times x_2=\tfrac12\times\tfrac12
\end{align*}
不在所有集合 $A$ 的交之中.

\par
\begin{figure}[H]
\centering
\includegraphics[width=0.42\textwidth]{fig37-1.png}
\caption{}\label{figure:ch05-fig37-1}
\end{figure}

你也许会说,``你选的不合适嘛! 如果在选定了 $x_1=\tfrac12$ 之后, 再选 $x_2=\tfrac23$, 那么就找到了交 $\bigcap_{A\in\mathcal{A}}A$ 的一个点.'' 然而, 因为在我们所设计的试探性的证明中, $x_1$ 和 $x_2$ 的选取是具有随意性的, 所以这种``不好''的选取便难以避免了.

那么, 怎样改进证明解决这个问题呢?

这就引出了证明的第二个想法: 也许我们可以扩大族 $\mathcal{A}$(当然还要保持有限交性质), 藉此来制约 $x_1$ 和 $x_2$ 的选取, 以确保我们做出一个``正确的''选择. 例如, 在前面的例子中, 把集族 $\mathcal{A}$ 扩大为以下椭圆形区域族 $\mathcal{D}$, 其中每一个区域的边界都是以 $p=\bigl(\tfrac13,\tfrac13\bigr)$ 为一个焦点, 而另一个焦点落在线段 $pq$ 上. 集族 $\mathcal{D}$ 如\cref{figure:ch05-fig37-2}所示. 这个新族 $\mathcal{D}$ 仍具有有限交性质, 但是, 如果你想在交
\begin{align*}
\bigcap_{D\in\mathcal{D}}\overline{\pi_1(D)}
\end{align*}
中选出一点 $x_1$, 则只能选 $x_1=\tfrac13$. 类似地, 只能选 $x_2=\tfrac13$. 而 $\tfrac13\times\tfrac13$ 必然属于每一个集合 $D$, 因此也属于每一个集合 $A$. 换句话说, 将集族 $\mathcal{A}$ 扩大成集族 $\mathcal{D}$, 就能使我们得出理想的选择.

\par
\begin{figure}[H]
\centering
\includegraphics[width=0.40\textwidth]{fig37-2.png}
\caption{}\label{figure:ch05-fig37-2}
\end{figure}

当然, 在这个例子中, 我们精心地选择了 $\mathcal{D}$, 使得证明能够进行下去, 但在一般情况下, 我们对于 $\mathcal{D}$ 的选择应该有些什么要求呢? 这里就产生了证明的第三个想法: 为什么不能简单地要求选择``尽可能大''的 $\mathcal{D}$ —— 以至于没有比它更大的族具有有限交性质了, 再来看这个 $\mathcal{D}$ 是否满足要求呢? 这个族 $\mathcal{D}$ 的存在性并非显而易见, 需要证明. 并且证明中我们要用到 Zorn 引理. 但在我们证明了 $\mathcal{D}$ 的存在性之后, 实际上也就证明了 $\mathcal{D}$ 已大到足以保证我们得出恰当选择的程度了.

最后还需说明的是, 在我们的讨论中, 对于集族 $\mathcal{A}$ 中的每一个元素都是闭集的假设是没有必要的. 因为即使 $A\in\mathcal{A}$ 是闭集, $\pi_1(A)$ 也未必是闭集, 因此我们只要通过取闭包来得到具有紧致性的闭集. 基于上述原因, 我们只需对 $X$ 的任意具有有限交性质的子集族, 证明它们的闭包的交是非空的. 事实上, 这一方法的确是简便易行的.

\begin{lemma}\label{lemma:ch05-maximal-fip-family}
设 $X$ 是一个集合. $\mathcal{A}$ 是 $X$ 的一个具有有限交性质的子集族. 则存在 $X$ 的一个子集族 $\mathcal{D}$, 使得 $\mathcal{D}$ 包含 $\mathcal{A}$, $\mathcal{D}$ 具有有限交性质, 并且每一个以 $\mathcal{D}$ 为真子集的 $X$ 的子集族都不再具有有限交性质.
\end{lemma}

对于满足引理要求的族 $\mathcal{D}$, 我们通常称它关于有限交性质是极大的.

\begin{proof}
下面, 我们就应用 Zorn 引理来构造 $\mathcal{D}$. 回忆一下, 给定一个定义了严格偏序关系的集合 $A$, 如果 $A$ 的任意全序子集都有上界, 那么 $A$ 便有一个极大元.

需说明的是, 将要应用 Zorn 引理的集合 $A$, 它既不是 $X$ 的子集, 也不是 $X$ 的一个子集族, 而是以集族为元素的集合. 我们把这样的集合称之为``超集'', 并用一个空心字母表示它. 于是综述现在所用的记号就是,
\begin{align*}
&c \text{ 是 } X \text{ 的一个元素},\\
&C \text{ 是 } X \text{ 的一个子集},\\
&\mathcal{C} \text{ 是 } X \text{ 的一个子集族},\\
&\mathfrak{C} \text{ 是 } X \text{ 的子集族构成的超集}.
\end{align*}

根据假设, 给定 $X$ 的一个具有有限交性质的子集族 $\mathcal{A}$, 用 $\mathbb{A}$ 表示 $X$ 的所有满足 $\mathcal{B}\supset\mathcal{A}$ 及有限交性质的子集族 $\mathcal{B}$ 组成的超集. 我们用真包含关系作为 $\mathbb{A}$ 上的一个严格偏序. 为了证明引理, 我们需要说明 $\mathbb{A}$ 有一个极大元 $\mathcal{D}$.

为了应用 Zorn 引理, 需要说明的是: 若 $\mathfrak{B}$ 是 $\mathbb{A}$ 的一个``子超集'', 并且在真包含关系下是全序的, 则 $\mathfrak{B}$ 在 $\mathbb{A}$ 中有一个上界. 实际上我们将要证明集族
\begin{align*}
\mathcal{C}=\bigcup_{\mathcal{B}\in\mathfrak{B}}\mathcal{B}
\end{align*}
是 $\mathbb{A}$ 中的一个元素, 就是所求的上界.

为了说明 $\mathcal{C}$ 是 $\mathbb{A}$ 的一个元素, 就须证明 $\mathcal{C}\supset\mathcal{A}$, 并且 $\mathcal{C}$ 具有有限交性质. 首先根据 $\mathfrak{B}$ 中的每一个元素都包含着 $\mathcal{A}$, 可见 $\mathcal{C}\supset\mathcal{A}$. 为了证明 $\mathcal{C}$ 具有有限交性质, 设 $C_1,\cdots,C_n$ 是 $\mathcal{C}$ 的一些元素. 因为 $\mathcal{C}$ 是 $\mathfrak{B}$ 中一些元素的并, 所以对于每一个 $i$, 存在 $\mathfrak{B}$ 的一个元素 $\mathcal{B}_i$, 使得 $C_i\in\mathcal{B}_i$. 现在超集 $\{\mathcal{B}_1,\cdots,\mathcal{B}_n\}$ 包含于 $\mathfrak{B}$, 因此, 它也是由真包含关系决定的全序集. 因为它是有限的, 所以有一个最大元, 也就是说存在一个指标 $k$, 使得对于 $i=1,\cdots,n$ 有 $\mathcal{B}_i\subset\mathcal{B}_k$. 从而, 集合 $C_1,\cdots,C_n$ 都属于 $\mathcal{B}_k$. 又因为 $\mathcal{B}_k$ 具有有限交性质, 所以集合 $C_1,\cdots,C_n$ 有非空交.
\end{proof}

\begin{lemma}\label{lemma:ch05-maximal-family-properties}
设 $X$ 是一个集合, $\mathcal{D}$ 是 $X$ 的一个子集族并且关于有限交性质是极大的. 则
\begin{enumerate}[(a)]
\item $\mathcal{D}$ 中元素的任意有限交仍属于 $\mathcal{D}$.
\item 若 $A$ 是 $X$ 的一个子集并且与 $\mathcal{D}$ 中的每一个元素都相交, 则 $A$ 属于 $\mathcal{D}$.
\end{enumerate}
\end{lemma}

\begin{proof}
(a) 设 $B$ 是 $\mathcal{D}$ 中有限多个元素的交. 将 $B$ 加到 $\mathcal{D}$ 中来定义一个族 $\mathcal{E}$, 即 $\mathcal{E}=\mathcal{D}\cup\{B\}$. 我们证明 $\mathcal{E}$ 满足有限交性质, 于是由 $\mathcal{D}$ 的极大性可见 $\mathcal{E}=\mathcal{D}$, 因此 $B\in\mathcal{D}$.

取 $\mathcal{E}$ 中有限多个元素. 如果其中没有一个是 $B$, 那么根据 $\mathcal{D}$ 具有有限交性质, 可见它们的交非空. 如果其中有一个是 $B$, 则它们的交形如
\begin{align*}
D_1\cap\cdots\cap D_m\cap B.
\end{align*}
因为 $B$ 等于 $\mathcal{D}$ 中元素的有限交, 所以这个交非空.

(b) 给定 $A$, 定义 $\mathcal{E}=\mathcal{D}\cup\{A\}$. 我们证明 $\mathcal{E}$ 具有有限交性质, 从而推出 $A$ 属于 $\mathcal{D}$. 取 $\mathcal{E}$ 中有限多个元素, 如果其中没有 $A$, 它们的交当然非空. 否则, 它们的交形如
\begin{align*}
D_1\cap\cdots\cap D_n\cap A.
\end{align*}
根据(a)可见 $D_1\cap\cdots\cap D_n$ 属于 $\mathcal{D}$, 因此根据假设, 上述交非空.
\end{proof}

\begin{theorem}[Tychonoff 定理]\label{theorem:ch05-tychonoff}
在积拓扑下, 紧致空间的任意积还是紧致空间.
\end{theorem}

\begin{proof}
设
\begin{align*}
X=\prod_{\alpha\in J}X_\alpha,
\end{align*}
其中每一个空间 $X_\alpha$ 是紧致的. 设 $\mathcal{A}$ 是 $X$ 的一个子集族, 具有有限交性质, 我们证明交
\begin{align*}
\bigcap_{A\in\mathcal{A}}\overline{A}
\end{align*}
非空, 从而推出 $X$ 的紧致性.

应用\cref{lemma:ch05-maximal-fip-family}, 选择 $X$ 的子集族 $\mathcal{D}$, 使得 $\mathcal{D}\supset\mathcal{A}$, 并且关于有限交性质是极大的. 这就足以证明 $\bigcap_{D\in\mathcal{D}}\overline{D}$ 非空.

给定 $\alpha\in J$, 令 $\pi_\alpha\colon X\to X_\alpha$ 为通常的投射. 考虑 $X_\alpha$ 的子集族
\begin{align*}
\{\pi_\alpha(D)\mid D\in\mathcal{D}\}.
\end{align*}
它具有有限交性质, 这是因为 $\mathcal{D}$ 具有有限交性质. 由 $X_\alpha$ 的紧致性, 可以对于每一个 $\alpha$, 选取一点 $x_\alpha\in X_\alpha$, 使得
\begin{align*}
x_\alpha\in\bigcap_{D\in\mathcal{D}}\overline{\pi_\alpha(D)}.
\end{align*}
令 $x$ 为 $X$ 的一个点 $(x_\alpha)_{\alpha\in J}$, 我们将证明对于每一个 $D\in\mathcal{D}$, $x\in\overline{D}$, 从而完成定理的证明.

首先证明, 如果 $\pi_\beta^{-1}(U_\beta)$ 是包含 $x$ 的任意子基元(就 $X$ 的积拓扑而言), 那么 $\pi_\beta^{-1}(U_\beta)$ 与 $\mathcal{D}$ 的任何成员都相交, 这里 $U_\beta$ 是 $x_\beta$ 在 $X_\beta$ 中的一个邻域. 因为根据定义 $x_\beta\in\overline{\pi_\beta(D)}$, $U_\beta$ 与 $\pi_\beta(D)$ 交于某一点 $\pi_\beta(y)$, 其中 $y\in D$. 于是 $y\in\pi_\beta^{-1}(U_\beta)\cap D$.

根据\cref{lemma:ch05-maximal-family-properties}的结论(b)可见包含 $x$ 的每一个子基元都属于 $\mathcal{D}$. 于是, 再由这个引理的结论(a)可见每一个包含 $x$ 的基元也属于 $\mathcal{D}$. 因为 $\mathcal{D}$ 具有有限交性质, 这就意味着每一个包含 $x$ 的基元与 $\mathcal{D}$ 的每一个元素都相交, 因此对于每一个 $D\in\mathcal{D}$, $x\in\overline{D}$.
\end{proof}

\subsection*{习题}

\begin{enumerate}
\item 设 $X$ 是一个空间, $\mathcal{D}$ 是 $X$ 的关于有限交性质的一个极大子集族,
\begin{enumerate}[(a)]
\item 证明: 对于每一个 $D\in\mathcal{D}$, $x\in\overline{D}$ 当且仅当 $x$ 的每一个邻域属于 $\mathcal{D}$, 试指出, 证明哪一个蕴涵关系时用到了 $\mathcal{D}$ 的极大性?
\item 设 $D\in\mathcal{D}$. 证明: 如果 $A\supset D$, 那么 $A\in\mathcal{D}$.
\item 证明: 若 $X$ 满足 $T_1$ 公理, 则最多有一个点属于 $\bigcap_{D\in\mathcal{D}}\overline{D}$.
\end{enumerate}
\item 若 $X$ 的一个子集族 $\mathcal{A}$ 中任意可数个元素的交非空, 则称 $\mathcal{A}$ 满足可数交性质(countable intersection property). 证明: $X$ 是 Lindelöf 空间当且仅当 $X$ 的每一个满足可数交性质的子集族 $\mathcal{A}$, 有
\begin{align*}
\bigcap_{A\in\mathcal{A}}\overline{A}\neq\varnothing.
\end{align*}
\item 考虑下述三个论断:
\begin{enumerate}[(i)]
\item 若 $X$ 是一个集合, $\mathcal{A}$ 是 $X$ 的一个满足可数交性质的子集族, 则存在 $X$ 的一个子集族 $\mathcal{D}$, 使得 $\mathcal{D}\supset\mathcal{A}$, 并且 $\mathcal{D}$ 关于可数交性质是极大的.
\item 设 $\mathcal{D}$ 关于可数交性质是极大的, 则 $\mathcal{D}$ 的元素的可数交仍然是 $\mathcal{D}$ 的元素, 并且如果 $A$ 是 $X$ 的一个子集, 它与 $\mathcal{D}$ 的每一个元素都相交, 那么必然有 $A\in\mathcal{D}$.
\item Lindelöf 空间的积仍然是 Lindelöf 空间.
\end{enumerate}
\begin{enumerate}[(a)]
\item 证明: (i)和(ii)蕴涵着(iii).
\item 证明(ii).
\item Lindelöf 空间的积未必是 Lindelöf 空间(见第 30 节), 因此(i)不成立. 试指出若将\cref{lemma:ch05-maximal-fip-family}的证明推广到可数交性质, 那么证明在什么地方行不通.
\end{enumerate}
\item 下面是另一个应用 Zorn 引理证明的定理. 以前曾经提到过: 设 $A$ 是一个空间, 若 $x,y\in A$, 不存在 $A$ 的分割 $A=C\cup D$($C$ 和 $D$ 是 $A$ 中两个无交的开集), 使得 $x\in C$, $y\in D$, 则称 $x$ 和 $y$ 在 $A$ 的同一个拟连通分支中.

\textbf{定理} 设 $X$ 是紧致的 Hausdorff 空间, 则 $x$ 和 $y$ 属于 $X$ 的同一拟连通分支当且仅当它们属于 $X$ 的同一分支.
\begin{enumerate}[(a)]
\item 令 $\mathcal{A}$ 为 $X$ 的闭子空间的族, 满足条件: 对于每一个 $A\in\mathcal{A}$, $x,y$ 在 $A$ 的同一拟连通分支中. 令 $\mathcal{B}$ 为 $\mathcal{A}$ 的一个子族, 在真包含关系下是全序集. 证明 $\mathcal{B}$ 的所有元素的交属于 $\mathcal{A}$. [提示: 与第 26 节的练习 11 比较.]
\item 证明 $\mathcal{A}$ 有极小元 $D$.
\item 证明 $D$ 是连通的.
\end{enumerate}
\item 下面是 Tychonoff 定理的另外一个证明, 它主要利用了良序定理而非 Zorn 引理. 先来证明管状引理的以下形式, 再证明定理.

\textbf{引理} 设 $\mathcal{A}$ 是 $X\times Y$ 积拓扑的基元素构成的族, 满足条件: $\mathcal{A}$ 没有有限子族覆盖 $X\times Y$. 若 $X$ 是紧致的, 则存在点 $x\in X$ 使得 $\mathcal{A}$ 的有限子族都不能够覆盖薄片 $\{x\}\times Y$.

\textbf{定理} 任意紧致空间的积在积拓扑下是紧致的.

证明: 设 $\{X_\alpha\}_{\alpha\in J}$ 是紧致空间的一个加标族. 令
\begin{align*}
X=\prod_{\alpha\in J}X_\alpha,
\end{align*}
$\pi_\alpha\colon X\to X_\alpha$ 为投射. 赋予 $J$ 一个良序, 使得有一个最大元.
\begin{enumerate}[(a)]
\item 令 $\beta\in J$. 对于任意 $i<\beta$ 假设已经给定 $p_i\in X_i$. 对于任意 $\alpha<\beta$, 定义 $X$ 的子空间 $Y_\alpha$ 为
\begin{align*}
Y_\alpha=\{x\mid \pi_i(x)=p_i,\text{ 对 }i\leqslant\alpha\}.
\end{align*}
注意, 若 $\alpha<\alpha'$, 则 $Y_\alpha\supset Y_{\alpha'}$. 证明: 若 $\mathcal{A}$ 是 $X$ 的基元素构成的一个有限族, 并且 $\mathcal{A}$ 覆盖
\begin{align*}
Z_\beta=\bigcap_{\alpha<\beta}Y_\alpha=\{x\mid \pi_i(x)=p_i,\,i<\beta\},
\end{align*}
则对于某一个 $\alpha<\beta$, $\mathcal{A}$ 覆盖 $Y_\alpha$. [提示: 若在 $J$ 中 $\beta$ 有一个紧接前元, 用 $\alpha$ 表示这个紧接前元. 否则, 对于每一个 $A\in\mathcal{A}$, 令 $J_A$ 为使得 $\pi_i(A)\neq X_i$ 的那些指标 $i<\beta$ 的集合. 集合 $J_A(A\in\mathcal{A})$ 的并是有限的. 令 $\alpha$ 为这个并的最大元.]
\item 设 $\mathcal{A}$ 是 $X$ 的基元素的族, 满足条件: $\mathcal{A}$ 没有有限子族覆盖 $X$. 证明对于所有 $i$, 可以选取点 $p_i\in X_i$, 使得对于每一个 $\alpha$, (a)中所定义的空间 $Y_\alpha$ 不能被 $\mathcal{A}$ 有限子族覆盖. 当 $\alpha$ 是 $J$ 的最大元时, 便有了矛盾. [提示: 若 $\alpha$ 是 $J$ 的最小元, 则应用前面的引理选择 $p_\alpha$. 若 $p_i$ 对于所有 $i<\beta$ 有定义, 由于根据(a)可见空间 $Z_\beta$ 不能被 $\mathcal{A}$ 的有限子族覆盖, 应用这个引理找出 $p_\beta$.]
\end{enumerate}
\end{enumerate}

\end{document}